%=====================================================
% Dados de identificação exigidos pelo Anexo I.
%=====================================================

\chapter*{IDENTIFICAÇÃO}

\begin{spacing}{1.15}

\noindent
\begin{tabular}{@{}p{0.26\linewidth}p{0.70\linewidth}@{}}
\toprule
\textbf{Estudante} & André Victor Xavier Pires \\[0.4em]
\textbf{Matrícula} & GRR20212735 \\[0.4em]
\textbf{Curso} & Engenharia Elétrica --- Setor de Tecnologia, Universidade
Federal do Paraná \\[0.4em]
\textbf{Ênfase} & Sistemas eletrônicos embarcados (curso noturno) \\[0.4em]
\textbf{Equipe} & Trabalho individual \\
\midrule
\textbf{Orientador} & Prof.\ Dr.\ Leandro dos Santos Coelho --- Departamento de
Engenharia Elétrica e Programa de Pós-Graduação em Engenharia Elétrica
(PPGEE), UFPR \\[0.4em]
\textbf{Coorientador} & Prof.\ Dr.\ Bruno Pohlot Ricobom \\[0.4em]
\textbf{\quad Vínculo} & Departamento de Engenharia Elétrica, Setor de
Tecnologia, Universidade Federal do Paraná \\[0.4em]
\textbf{\quad Área de atuação} & Instrumentação eletrônica e medição de campo
próximo aplicada a compatibilidade eletromagnética \\[0.4em]
\textbf{\quad Contribuição} & Instrumentação e medição: concepção das placas de
teste e condução da validação experimental por varredura de campo próximo
(Seção~\ref{subsec:validacao}) \\
\midrule
\textbf{Disciplinas} & TCC I (período letivo 2026/2) e TCC II (período letivo
2027/1) \\
\bottomrule
\end{tabular}

\vspace{2em}

\noindent
A descrição do projeto que segue reproduz, na ordem, os itens enumerados no
Anexo I das normas do curso: o item 1 (título) consta da folha de rosto; os
itens 2 a 9 são as seções do Capítulo~\ref{cap:descricao}, cujos títulos
repetem literalmente a designação do Anexo; e o item 10 (bibliografia a ser
utilizada) corresponde a REFERÊNCIAS. O sumário serve, portanto, como o próprio
quadro de conferência do Anexo I.

\end{spacing}

\vfill
\cleardoublepage

%=====================================================
